\section{Himpunan, Relasi, Fungsi, dan Induksi Matematika}
\label{sec:ch02-himpunan}

Teori kompilasi dibangun di atas bahasa matematika yang ringkas dan
presisi. Sebelum membahas token, tata bahasa, atau otomata, kita perlu
menyepakati notasi dasar. Bagian ini merangkum konsep himpunan, relasi,
fungsi, dan induksi matematika yang akan dipakai berulang kali.

\subsection{Himpunan dan Operasi Himpunan}

\begin{definisi}[Himpunan dan keanggotaan]
Himpunan adalah kumpulan objek yang terdefinisi dengan baik. Jika \(x\)
merupakan anggota himpunan \(A\), kita menulis \(x \in A\); jika tidak,
\(x \notin A\). Himpunan yang tidak memiliki anggota disebut himpunan
kosong dan dilambangkan \(\varnothing\).
\end{definisi}

Himpunan dapat dinyatakan dengan mendaftar anggotanya, misalnya
\(A = \{0,1,2\}\), atau dengan aturan keanggotaan, misalnya
\(B = \{n \in \mathbb{N} \mid n \text{ genap}\}\). Himpunan \(A\)
merupakan himpunan bagian dari \(B\), ditulis \(A \subseteq B\), jika
setiap anggota \(A\) juga anggota \(B\).

Operasi dasar himpunan adalah gabungan \(A \cup B\), irisan
\(A \cap B\), selisih \(A \setminus B\), dan komplemen \(A^c\) terhadap
semesta pembicaraan. Himpunan kuasa \(\mathcal{P}(A)\) adalah himpunan
seluruh himpunan bagian \(A\). Jika \(A\) memiliki \(n\) anggota, maka
\(\mathcal{P}(A)\) memiliki \(2^n\) anggota.

Produk Cartesian \(A \times B\) adalah himpunan seluruh pasangan
terurut \((a,b)\) dengan \(a \in A\) dan \(b \in B\). Konsep ini menjadi
dasar untuk mendefinisikan relasi dan fungsi.

\subsection{Relasi dan Fungsi}

\begin{definisi}[Relasi dan fungsi]
Relasi dari \(A\) ke \(B\) adalah himpunan bagian dari \(A \times B\).
Fungsi \(f : A \to B\) adalah relasi khusus yang memasangkan setiap
\(a \in A\) dengan tepat satu \(b \in B\). Himpunan \(A\) disebut
domain, dan himpunan nilai \(f(a)\) disebut range.
\end{definisi}

Fungsi \(f\) bersifat \emph{injektif} jika \(f(a_1) = f(a_2)\)
mengakibatkan \(a_1 = a_2\), \emph{surjektif} jika setiap anggota \(B\)
merupakan nilai \(f\), dan \emph{bijektif} jika kedua sifat tersebut
dimiliki. Dalam teori bahasa, fungsi sering dipakai untuk memetakan
string ke string lain, misalnya fungsi pembalik \(w \mapsto w^R\).

Relasi ekuivalensi adalah relasi yang bersifat reflektif, simetris, dan
transitif. Relasi ini mempartisi himpunan menjadi kelas-kelas
ekuivalensi. Konsep partisi akan muncul kembali saat minimalisasi DFA:
state yang berperilaku sama dikelompokkan ke dalam kelas ekuivalensi.

\subsection{Induksi Matematika}

\begin{definisi}[Prinsip induksi sederhana]
Untuk membuktikan bahwa pernyataan \(P(n)\) benar untuk setiap bilangan
asli \(n \ge n_0\), cukup dibuktikan:
\begin{enumerate}[label=(\roman*),leftmargin=*]
  \item \textbf{basis}: \(P(n_0)\) benar;
  \item \textbf{langkah induksi}: jika \(P(k)\) benar untuk suatu
        \(k \ge n_0\), maka \(P(k+1)\) benar.
\end{enumerate}
\end{definisi}

Induksi kuat mengizinkan asumsi bahwa \(P(n_0), P(n_0+1), \ldots, P(k)\)
semuanya benar untuk membuktikan \(P(k+1)\). Induksi struktural adalah
varian yang dipakai untuk membuktikan sifat string, pohon, atau
ekspresi yang didefinisikan secara rekursif.

Sebagai contoh, kita buktikan bahwa panjang konkatenasi memenuhi
\(|uv| = |u| + |v|\) untuk semua string \(u\) dan \(v\). Basis:
jika \(v = \varepsilon\), maka \(uv = u\) dan \(|uv| = |u| =
|u| + 0 = |u| + |v|\). Langkah induksi: misalkan sifat berlaku untuk
\(v\). Untuk \(v' = va\) dengan \(a\) sebuah simbol, berlaku
\(uv' = (uv)a\), sehingga
\[
|uv'| = |uv| + 1 = |u| + |v| + 1 = |u| + |v'|.
\]
Dengan induksi struktural, sifat terbukti untuk semua string \(v\).

Induksi juga dipakai untuk membuktikan bahwa sebuah tata bahasa
membangkitkan bahasa tertentu. Misalnya, untuk membuktikan bahwa
\(G\) membangkitkan \(L(G) = \{a^n b^n \mid n \ge 0\}\), kita
membuktikan dua arah: setiap string yang diturunkan memiliki bentuk
tersebut, dan setiap string berbentuk tersebut dapat diturunkan.
